% !TEX root = ../../tir_monograph_v12.tex
% V12_SOURCE_MAP:
% - TIR/foundations/TIR_KAPPA_FLAVOUR_MIXING_NORMALIZATION_V0_1.md
% - TIR/monograph/appendices/appC_su3_primer.tex
% - TIR/monograph/chapters/ch16_pmns_mixing.tex
% - TIR/monograph/chapters/ch18_ckm_matrix.tex
\chapter[Three-Flavour Carrier]{Three-Flavour Carrier and $SU(3)_F$}
\label{ch:v12_three_flavour_su3}

\paragraph{Established context.}
For a complex three-component carrier, unitary basis changes are described by
$U(3)$, while determinant-one unitary transformations form $SU(3)$.  The Lie
algebra dimension
\[
\dim_{\mathbb R}\mathfrak{su}(N)=N^2-1
\]
is standard.  TIR uses this group-theoretic structure as the family-mixing
carrier for its three-flavour branch.

\vTwelveStatus{B}{--}{PASS}

\section{Three-flavour carrier}
The family branch is represented by
\begin{equation}
\boxed{V_F\cong\mathbb C^3.}
\label{eq:v12-flavour-carrier}
\end{equation}
Choose an ordered flavour basis
\[
\{|f_1\rangle,|f_2\rangle,|f_3\rangle\}.
\]
A normalized family state is
\[
|\psi_F\rangle
=\sum_{i=1}^{3}c_i|f_i\rangle,
\qquad
\sum_{i=1}^{3}|c_i|^2=1.
\]
Full determinant-one family mixing acts through
\begin{equation}
\boxed{U_F\in SU(3)_F.}
\label{eq:v12-su3f-carrier}
\end{equation}
The carrier dimension therefore fixes the flavour multiplicity
\begin{equation}
\boxed{N_F=3.}
\label{eq:v12-flavour-multiplicity}
\end{equation}

\section{Eight independent mixing directions}
The associated mixing algebra is
\[
\mathfrak{su}(3)_F,
\]
with
\begin{equation}
\boxed{
\dim_{\mathbb R}\mathfrak{su}(3)_F
=3^2-1
=8.
}
\label{eq:v12-su3f-dimension}
\end{equation}
If $\{T_a\}_{a=1}^{8}$ is any basis of generators, a local family
transformation may be written
\[
U_F(\boldsymbol\theta)
=\exp\!\left(i\sum_{a=1}^{8}\theta_aT_a\right).
\]
The number eight is basis-independent: it is the dimension of the Lie algebra,
not a property of a particular generator convention.

The TIR symmetric-pair crosswalk supplies the decomposition
\begin{equation}
\mathfrak{su}(3)_F
=\mathfrak{so}(3)\oplus\mathfrak p,
\qquad
8=3+5,
\label{eq:v12-su3-symmetric-pair}
\end{equation}
which gives an independent carrier-level accounting of the same eight real
mixing directions.

\section{Generator--flavour incidence set}
For later normalization bookkeeping define
\begin{equation}
\mathcal C_{\rm mix}
=\{1,\ldots,8\}_{\rm generator}
\times\{1,2,3\}_{\rm flavour}.
\label{eq:v12-mixing-incidence-set}
\end{equation}
Its cardinality is
\begin{equation}
|\mathcal C_{\rm mix}|
=8\cdot3
=24.
\label{eq:v12-mixing-incidence-cardinality}
\end{equation}
At this stage the integer $24$ is only the exact incidence count generated by
the already specified carrier and algebra.  Chapter~\ref{ch:v12_kappa_normalization}
adds the half-turn phase measure and the binary-information numerator to obtain
the TIR normalization.

\section{Downstream mixing matrices}
The PMNS and CKM constructions are treated in
Chapter~\ref{ch:v12_ckm_pmns}.  Their numerical angles, phases and empirical
residuals are downstream sector assignments.  The carrier statements
\eqref{eq:v12-flavour-carrier}--\eqref{eq:v12-su3f-dimension} are the structural
parents shared by those sectors and by the normalization chapter that follows.
